\documentclass[conference]{IEEEtran}
\usepackage[T1]{fontenc}
\usepackage[utf8]{inputenc}
\usepackage{amsmath,amssymb}
\usepackage{graphicx}
\usepackage{booktabs}
\usepackage{cite}
\usepackage{url}

\newcommand{\tabref}[1]{Table~\ref{#1}}
\newcommand{\figref}[1]{Fig.~\ref{#1}}
\newcommand{\secref}[1]{Section~\ref{#1}}
\newcommand{\Ldec}{\mathcal{L}_\mathrm{dec}}
\newcommand{\Lenc}{\mathcal{L}_\mathrm{enc}}

% Keep captions, footnotes, and references at the existing 9 pt minimum.
\renewcommand{\footnotesize}{\small}

\begin{document}

\title{Mixing Supervision and Latent Arithmetic\\in Audio Autoencoders}

\author{\IEEEauthorblockN{Nurdaulet Akhanov}
\IEEEauthorblockA{Mohamed bin Zayed University of Artificial Intelligence\\
Abu Dhabi, United Arab Emirates\\
nurdaulet.akhanov@mbzuai.ac.ae}}

\maketitle

\begin{abstract}
Evaluations of latent arithmetic can conflate waveform fidelity, decoder
agreement, and the choice of latent origin. We examine these effects through
training comparisons in audio autoencoders trained with a generative
adversarial network (GAN) objective and independent held-out evaluation.
Decode-mixing supervision improves the scale-invariant signal-to-distortion
ratio (SI-SDR) of waveform mixing by
$1.77$\,dB at $7.66\times$ compression; a combined recipe improves it by
$1.39$\,dB at $15.3\times$. Adversarial supervision on the mixed path
increases decoder agreement while reducing mixing fidelity. For source
subtraction, recentering at each model's encoding of silence removes
$87$--$95$\% of the raw between-model advantage on 240 held-out recordings,
while preserving reconstruction and convex mixing exactly. These findings
separate waveform mixing gains from decoder agreement and origin-dependent
subtraction effects. Results are conditional on single training runs; the
independent evaluation uses frozen checkpoints selected on the original test
split.
\end{abstract}

\begin{IEEEkeywords}
audio autoencoders, latent arithmetic, mixing supervision,
representation learning, evaluation methodology
\end{IEEEkeywords}

\section{Introduction}\label{sec:intro}
Mixing audio is a linear operation on waveforms. Performing the same
operation on latent representations would provide a simple interface for
composition and editing. For an encoder $f$ and decoder $g$, the desired
property is
\begin{equation}
  g\!\left(\alpha f(x_1)+(1{-}\alpha)f(x_2)\right)
  \approx \alpha x_1+(1{-}\alpha)x_2,
  \label{eq:equiv}
\end{equation}
where $\alpha\in[0,1]$. A related operation removes a known stem by
subtracting its encoding from that of a mixture. Such arithmetic could also
support composition from independently generated latents in audio diffusion
systems \cite{liu2023audioldm,evans2024stable}. Reconstruction objectives
alone, however, do not enforce these algebraic relationships.

Explicit mixing supervision addresses this gap by training a decoder to
recover a waveform mixture from a combination of separately encoded inputs
\cite{torres2026linearity}. Assessing its benefit requires care. Comparing
the decoded combination with the true mixture measures waveform fidelity;
comparing it with another output of the same decoder measures agreement
between decoding paths. Source subtraction introduces a further dependency:
its output can change when the latent origin is shifted, even if
reconstruction and convex mixing remain identical.

We investigate these distinctions in waveform GAN autoencoders at two
compression rates. Our contribution is an empirical study of an existing
supervision principle and of the measurements used to assess it:
\begin{enumerate}
  \item We compare decode-mixing supervision, encoder-side mixing
  supervision, and decoder-only adaptation from a shared checkpoint.
  Decode-mixing yields the strongest waveform mixing result, and
  decoder-only adaptation retains most of the gain.
  \item We show that adversarial supervision on the mixed path improves
  decoder agreement while reducing waveform mixing fidelity. We also
  examine how decoder sampling and latent-error normalization affect
  evaluation.
  \item We establish that recentering at encoded silence preserves
  reconstruction and all unit-sum latent combinations, then show that it
  removes most of the apparent subtraction advantage of mixing supervision.
\end{enumerate}
\section{Related Work}\label{sec:related}
\textbf{Audio representations.} Neural codecs such as DAC
\cite{kumar2024dac}, EnCodec \cite{defossez2023encodec}, and SoundStream
\cite{zeghidour2022soundstream} use convolutional encoders, learned decoders,
and adversarial reconstruction objectives. Interpolating their quantized
latents generally leaves the discrete set of codes.
Continuous representations, including RAVE \cite{caillon2021rave},
avoid this restriction but do not explicitly enforce mixing equivariance.
Disentanglement objectives such as $\beta$-VAE \cite{higgins2017betavae}
likewise target a different property. Relating waveform operations to latent
operations follows the broader study of structure in representations
\cite{higgins2018towards}; here, convex mixing is an affine combination
rather than a group action.

\textbf{Mixing supervision.} Music2Latent (M2L) \cite{pasini2024music2latent}
learns a continuous audio representation with a consistency decoder. Torres
et al.~\cite{torres2026linearity} explicitly encourage linearity by decoding
the average of separate encodings against the corresponding waveform
mixture, a principle related to mixup \cite{zhang2018mixup}. We apply this
principle to a waveform GAN autoencoder, replacing the consistency loss with
spectral and waveform losses and sampling $\alpha$ uniformly instead of
fixing it at $0.5$. Prior additivity and separation
scores compare two decoded outputs \cite{torres2026linearity}; we examine
that reference choice alongside ground-truth waveform evaluation.

\section{Method and Evaluation}\label{sec:method}

\subsection{Mixing objectives}\label{sec:loss}
We pair batch elements $(x_i,x_j)$ using a permutation with no fixed points
and draw $\alpha\sim U[0,1]$. Define the waveform mixture
$\bar{x}=\alpha x_i+(1{-}\alpha)x_j$ and latent interpolation
$\bar{z}=\alpha f(x_i)+(1{-}\alpha)f(x_j)$. The decode-mixing objective is
\begin{equation}
  \Ldec=\mathcal{L}_\mathrm{recon}\!\left(g(\bar{z}),\bar{x}\right)
  +\mathcal{L}_{\mathrm{wave}\text{-}\ell_1}\!\left(g(\bar{z}),\bar{x}\right),
  \label{eq:ldec}
\end{equation}
where $\mathcal{L}_\mathrm{recon}$ combines equally weighted
multi-resolution STFT and mel losses. The waveform term is the unweighted
mean absolute error over samples and batch elements. It is applied only to
the mixed path. Thus, this recipe introduces three coupled changes: an
interpolated latent, a mixture target, and an additional waveform loss.
The comparison with the baseline measures their combined effect.

Gradients from $\Ldec$ pass through both $g$ and $f$. We separately test an
encoder-side objective,
$\Lenc=\lVert f(\bar{x})-\bar{z}\rVert^2$, a decode-mixing variant with
$f$ frozen, and a variant that also applies adversarial supervision to the
mixed path (``$+$disc''). The training objective is
$\mathcal{L}=\mathcal{L}_\mathrm{recon}+\mathcal{L}_\mathrm{GAN}
+\lambda\Ldec+\gamma\Lenc+\lambda_{\ell_2}\lVert z\rVert^2$.
The baseline uses $\lambda=\gamma=0$; decode-mixing uses $\lambda=0.5$ and
$\gamma=0$; encoder-side mixing uses $\lambda=0$ and
$\gamma\in\{5,10,20\}$. The $+$disc variant is also the combined recipe
used at $15.3\times$ compression.

\subsection{Models, training, and data}\label{sec:setup}
Our primary model combines a DAC-style 1-D convolutional encoder, a
HiFi-GAN decoder \cite{kong2020hifigan}, and multi-scale STFT and multi-period
discriminators. It processes 1-second mono chunks at 44.1\,kHz and is trained
on approximately 950 hours from FMA-large, MAESTRO, and MUSDB18-HQ
\cite{rafii2017musdb18}.\footnote{Code, configurations, pretrained weights,
evaluation records, and audio examples:
\url{https://github.com/nurdauletakhanov/musicgen}}
We study $7.66\times$ compression with latent dimension $d=128$ and
$15.3\times$ with $d=64$. At $7.66\times$, each configuration is fine-tuned
for 25k steps from the same converged checkpoint. At $15.3\times$, the
baseline and combined recipe are trained for 250k steps with matched
configurations but different initial weights.

All runs use $\lambda_{\ell_2}=10^{-3}$, adversarial weight $0.5$, and
feature-matching weight $2$. We use Adam with learning rate $10^{-4}$
($5{\times}10^{-5}$ for the discriminator), $\beta=(0.8,0.99)$, weight decay
$10^{-3}$, gradient clipping at $1.0$, batch size $32$, two gradient
accumulation steps, and $500$ warmup steps.

Checkpoints were selected by the lowest validation loss, but the validation
loader used the original test split. The selected checkpoint is the final
step for every fine-tune and step 240k for both from-scratch runs. We
therefore evaluate the five frozen GAN checkpoints on 240 MoisesDB
recordings \cite{pereira2023moisesdb} excluded from training and selection.
The protocol was registered before downloading this corpus; the released
\texttt{HOLDOUT\_PROTOCOL.md} records checkpoint hashes verified at evaluation
time. Held-out mixing results use a single set of $14{,}064$ paired units,
since changing the batching changes the pairs. We report 95\% intervals by
resampling recordings. A pair from distinct recordings receives bootstrap
weight $m_i m_j$, where $m_i$ is a recording's multiplicity; a pair from the
same recording receives weight $m_i$. About $80\%$ of pairs are within a
recording. These intervals quantify evaluation uncertainty for fixed
checkpoints.

\subsection{Waveform and latent metrics}\label{sec:metrics}
We use SI-SDR \cite{leroux2019sdr} for waveform comparisons. For an estimate
$\hat{x}$ and reference $x$, each with its temporal mean removed, define
\begin{equation}
  s(\hat{x},x)=10\log_{10}
  \frac{\lVert a x\rVert_2^2}{\lVert\hat{x}-a x\rVert_2^2},
  \qquad a=\frac{\langle\hat{x},x\rangle}{\lVert x\rVert_2^2}.
  \label{eq:sisdr}
\end{equation}
The gain $a$ accounts for level mismatch. We stabilize energy terms with
$\epsilon=10^{-8}$ and use the same score for reconstruction and mixing:
\begin{align}
  \mathrm{SI\text{-}SDR}_\mathrm{rec}
    &=s(g(f(x)),x), \label{eq:recmetric}\\
  \mathrm{SI\text{-}SDR}_\mathrm{lin}^\mathrm{gt}
    &=s(g(\bar{z}),\bar{x}), \label{eq:gtmetric}\\
  \mathrm{SI\text{-}SDR}_\mathrm{lin}
    &=s(g(\bar{z}),g(f(\bar{x}))). \label{eq:consmetric}
\end{align}
Mixing is evaluated at $\alpha=0.5$ unless stated otherwise.
Equation~\ref{eq:gtmetric} measures fidelity to the true mixture;
Eq.~\ref{eq:consmetric} measures agreement with its reconstruction. We also report
MixRate, the ratio of the loss in Eq.~\ref{eq:ldec} for $g(\bar{z})$ to that
for $g(f(\bar{x}))$, both against $\bar{x}$. A value below one means the
interpolated route has lower loss than direct reconstruction. Ground-truth
SI-SDR is our primary mixing measure; MixRate supplies a complementary
comparison with each model's reconstruction path.

For the latent residual $r=\bar{z}-f(\bar{x})$, the conventional normalized
error is $\ell_\mathrm{lat}=\lVert r\rVert^2/\lVert f(\bar{x})\rVert^2$.
We examine alternative normalizations in \secref{sec:training}; latent
statistics are computed per sample and averaged over evaluation units.
Because SI-SDR fits a gain before scoring, we assess that gain separately on
3{,}000 samples. FAD uses music-finetuned LAION-CLAP embeddings of
non-overlapping 10\,s windows at 48\,kHz, so its values are not directly
comparable with VGGish-based FAD. Paired stochastic decodes use shared noise,
except in the noise-control experiment in \secref{sec:consistency}.

\subsection{Recentering the latent origin}\label{sec:prop}
Let $c=f(0)$ be the encoding of silence and define
$f_c(x)=f(x)-c$ and $g_c(z)=g(z+c)$.
\textbf{Proposition.} For arbitrary $f$ and $g$, this change of coordinates
preserves reconstruction and every latent combination whose coefficients
sum to one:
$g_c(f_c(x))=g(f(x))$ and
$g_c(\sum_i a_i f_c(x_i))=g(\sum_i a_i f(x_i))$ when $\sum_i a_i=1$.
Subtraction instead becomes
\begin{equation}
  g_c\!\left(f_c(x)-f_c(y)\right)
  =g\!\left(f(x)-f(y)+f(0)\right).
  \label{eq:sub}
\end{equation}
\emph{Proof.} Reconstruction cancels the subtracted offset when $g_c$ adds
it back. For a unit-sum combination,
$\sum_i a_i f_c(x_i)=\sum_i a_i f(x_i)-c$, so the same cancellation holds.
In a difference, the two encoder offsets cancel each other, leaving the
offset added by $g_c$. $\square$

The identities hold exactly for nonlinear networks and require no
retraining. They show that subtraction scores depend on the chosen origin
even when reconstruction and convex mixing are fixed. If $f$ is affine,
$f(x)-f(y)+f(0)=f(x-y)$, so corrected subtraction returns the model's own
reconstruction of the residual. We use that reconstruction as a reference
and call the SI-SDR shortfall the \emph{linearity tax}. For nonlinear
encoders, any improvement from recentering is empirical. The proposition
concerns frozen models: it does not make training invariant to the origin,
and the latent $\ell_2$ penalty can influence $f(0)$.

\section{Results}\label{sec:results}

\subsection{Waveform mixing improves on independent audio}\label{sec:mixing}
On MoisesDB, decode-mixing improves ground-truth mixing SI-SDR by
$1.77$\,dB $[1.73,1.81]$ at $7.66\times$ compression. The combined recipe
improves it by $1.39$\,dB $[1.33,1.45]$ at $15.3\times$. These findings
support the original test-set comparison in \tabref{tab:ablation}, where
$\Ldec$ increases SI-SDR$_\mathrm{lin}^\mathrm{gt}$ from $8.0$ to
$10.2$\,dB and is the only configuration with MixRate below one ($0.97$).
Across these fine-tunes, reconstruction SI-SDR ranges from $11.2$ to
$11.4$\,dB and FAD from $0.044$ to $0.046$.

\begin{table}[tb]
  \centering\small
  \setlength{\tabcolsep}{3.0pt}
  % AUTO-GENERATED by make_tables.py � do not edit by hand (regenerate via make_tables.py).
\begin{tabular}{@{}lccccc@{}}
\toprule
Loss & SDR$_\mathrm{rec}$ & SDR$_\mathrm{lin}$ & SDR$_\mathrm{lin}^\mathrm{gt}$ & $\ell_\mathrm{lat}$ & MixRate \\
\midrule
    baseline (no mix) & +11.3 & +12.0 & +8.0 & 0.068 & 1.203 \\
    $\mathcal{L}_\mathrm{dec}$ & +11.4 & +12.1 & +10.2 & 0.052 & 0.972 \\
    $\mathcal{L}_\mathrm{dec}$ + disc-on-mix & +11.4 & +16.0 & +9.9 & 0.050 & 1.052 \\
    $\mathcal{L}_\mathrm{enc}$ ($\gamma{=}5$) & +11.3 & +12.8 & +8.4 & 0.044 & 1.175 \\
    $\mathcal{L}_\mathrm{enc}$ ($\gamma{=}10$) & +11.3 & +13.1 & +8.5 & 0.039 & 1.165 \\
    $\mathcal{L}_\mathrm{enc}$ ($\gamma{=}20$) & +11.3 & +13.5 & +8.7 & 0.034 & 1.153 \\
    $\mathcal{L}_\mathrm{dec}$ frozen-enc & +11.2 & +11.7 & +9.5 & 0.072 & 1.056 \\
\bottomrule
\end{tabular}

  \caption{GAN fine-tuning from a shared checkpoint ($7.66\times$, full
  original test set, $\alpha=0.5$; one run per configuration). SI-SDR
  subscripts follow Eqs.~\ref{eq:recmetric}--\ref{eq:consmetric}.}
  \label{tab:ablation}
\end{table}

The held-out gain largely reflects a smaller penalty for decoding
interpolated latents. Define this penalty as the ground-truth SI-SDR of
direct reconstruction $g(f(\bar{x}))$ minus that of $g(\bar{z})$, evaluated
on the same units. At $7.66\times$, direct reconstruction improves by
$0.04$\,dB $[0.03,0.05]$, while the penalty falls by $1.73$\,dB
$[1.70,1.77]$, from $1.80$ to $0.07$\,dB. These changes account for the
$1.77$\,dB mixing gain. At $15.3\times$, direct reconstruction changes by
$-0.14$\,dB $[-0.17,-0.10]$, while the penalty falls by $1.52$\,dB
$[1.48,1.58]$. This decomposition describes where the measured gain occurs;
it does not isolate a training mechanism. In a separate five-coefficient
sweep on the original data, $\Ldec$ outperforms its matched control at every
$\alpha$ and reduces the SI-SDR range from $3.2$ to $1.9$\,dB.

\subsection{Decoder adaptation and latent-error normalization}\label{sec:training}
The encoder-side objective produces the lowest conventional latent error
($0.034$--$0.044$) but smaller waveform mixing gains: SI-SDR improves from
$8.0$ to $8.4$--$8.7$\,dB, compared with $10.2$\,dB under $\Ldec$.
With the encoder frozen, decode-mixing still reaches $9.5$\,dB. Thus,
decoder-only adaptation retains most of the gain, while joint adaptation
performs best among the tested configurations. Whether decoder adaptation
is necessary remains untested because there is no decoder-frozen counterpart.

Latent error requires a separate qualification. Translating the encoder
outputs by a constant and compensating in the decoder preserves $r$ but
changes $\lVert f(\bar{x})\rVert^2$. A sufficiently large translation can
therefore make $\ell_\mathrm{lat}$ arbitrarily small without changing
waveform mixing. \tabref{tab:latnorm} compares this ratio with three
translation-invariant alternatives. Let $D$ be the number of latent elements
and define $N=\lVert r\rVert^2/D$,
$S=\lVert f(x_1)-f(x_2)\rVert^2/D$, and
$O=\lVert f(\bar{x})-f(0)\rVert^2/D$.
Both $N$ and $N/S$ favour mixing supervision in all three comparisons.
However, $N/O$ favours the controls at $7.66\times$ and shows no resolved
difference at $15.3\times$.

\begin{table}[tb]
  \centering\small
  \setlength{\tabcolsep}{3.0pt}
  % AUTO-GENERATED by make_tables.py — do not edit by hand.
\begin{tabular}{@{}lccc@{}}
\toprule
normalization of $r=\bar{z}-f(\bar{x})$ & $\Ldec$ & $\Ldec{+}$disc & combined \\
\midrule
    $\lVert r\rVert^2/\lVert f(\bar{x})\rVert^2$ \emph{(transl.-sens.)} & -29.9$\pm$1.7 & -36.0$\pm$1.9 & -45.1$\pm$2.8 \\
    $\lVert r\rVert^2/D$ & -2.3$\pm$0.1 & -2.6$\pm$0.2 & -3.0$\pm$0.4 \\
    $\lVert r\rVert^2/\lVert f(x_1){-}f(x_2)\rVert^2$ & -16.7$\pm$0.8 & -20.3$\pm$0.9 & -30.1$\pm$1.4 \\
    $\lVert r\rVert^2/\lVert f(\bar{x}){-}f(0)\rVert^2$ & +13.4$\pm$1.0 & +8.3$\pm$1.2 & -0.2$\pm$2.1$^{\dagger}$ \\
\bottomrule
\end{tabular}

  \caption{Held-out latent-error differences. Each recipe is compared with
  its own no-mixing control: the first two columns at $7.66\times$, the
  combined recipe at $15.3\times$. Values are treatment minus control
  ($\times10^{-3}$; negative favours treatment), $\pm$ half the 95\%
  recording-bootstrap interval. $\dagger$ marks the interval containing zero.}
  \label{tab:latnorm}
\end{table}

The ratios of treatment to control means help explain this disagreement.
Mean $N$ falls to $0.64$--$0.67$ of the control value. Mean $S$ stays within
$4\%$ of the control at $7.66\times$ and rises to $1.37$ times the control
at $15.3\times$; mean $O$ falls to $0.53$--$0.61$ of the control. Thus, the
mean squared residual decreases by about a third, while mixture latents move
$22$--$27\%$ closer to encoded silence in RMS distance. These ratios of
means do not determine the mean of each per-unit ratio, but the larger
reduction in $O$ is consistent with the reversal in $N/O$. Both $N/S$ and
$N/O$ also remain unchanged under a uniform latent rescaling with a
compensating decoder change; $N$ alone does not. Neither translation nor
scale invariance selects between the two normalized measures. We therefore
restrict our claim to improved waveform mixing, without inferring a general
improvement in encoder geometry.

\subsection{Decoder agreement can increase as fidelity decreases}\label{sec:consistency}
Adding adversarial supervision to the mixed path increases
decode-versus-decode SI-SDR from $12.1$ to $16.0$\,dB on the original test
set, while ground-truth mixing SI-SDR falls from $10.2$ to $9.9$\,dB and
MixRate rises from $0.97$ to $1.05$. The held-out comparison shows the same
% The 14,064-unit run in latent_invariance/mixing_stats.json gives
% +3.9847723926 dB; +4.06 in the previous draft came from the older run.
pattern: decoder agreement increases by $3.98$\,dB while ground-truth mixing
changes by $-0.49$\,dB $[-0.51,-0.48]$. The latent residual is also slightly
smaller ($0.0045$ versus $0.0047$ per element), despite a larger waveform
interpolation penalty ($0.52$ versus $0.07$\,dB). Agreement between decoding
paths, latent residual size, and fidelity to the waveform therefore rank
these checkpoints differently. The mechanism remains unestablished, and
mixed-path adversarial supervision without $\Ldec$ has not been tested.

Stochastic decoding introduces another dependency. For the pretrained
Music2Latent checkpoint, decoding the same latent twice gives identical outputs under shared
noise but an SI-SDR of $-1.8$\,dB under independent noise, without any latent
arithmetic. For the base and mix-trained M2L checkpoints on identical pairs,
independent noise gives decoder-agreement scores of $-9.1$ and
$-7.3$\,dB; shared noise gives $3.6$ and $5.9$\,dB. Changing only the noise
coupling therefore shifts the scores by approximately $13$\,dB. Prior work
\cite{torres2026linearity} does not specify this detail sufficiently for us
to determine which protocol it used.

\subsection{Recentering removes most of the subtraction advantage}\label{sec:origin}
We evaluate oracle stem removal: the removed stem is known and encoded
separately. Raw subtraction decodes
$f(x_\mathrm{mix})-f(x_\mathrm{stem})$; the recentered version adds $f(0)$
as in Eq.~\ref{eq:sub}. Both are scored against the true residual. The
MUSDB18 evaluation uses 5{,}880 units from 49 of the 50 test tracks; one track
is excluded because the sum of its stems matches its mixture at only
$18.5$\,dB SI-SDR.

\begin{table}[tb]
  \centering\small
  \setlength{\tabcolsep}{2.2pt}
  % AUTO-GENERATED by make_tables.py — do not edit by hand (regenerate via make_tables.py).
\begin{tabular}{@{}lcccccc@{}}
\toprule
Model & drums & bass & vocals & other & all & gap$\downarrow$ \\
\midrule
    7.66$\times$ no mix & +4.2 & +0.4 & +5.2 & +3.9 & +3.4 & 5.2 \\
    \quad $+f(0)$ & +8.8 & +6.8 & +9.6 & +8.0 & +8.3 & 0.4 \\
    7.66$\times$ $+\mathcal{L}_\mathrm{dec}$ & +5.7 & +4.9 & +7.2 & +5.1 & +5.7 & 3.0 \\
    \quad $+f(0)$ & +8.5 & +6.9 & +9.4 & +8.3 & +8.3 & 0.5 \\
    7.66$\times$ $+\mathcal{L}_\mathrm{dec}$+disc & +6.2 & +5.4 & +7.5 & +5.0 & +6.0 & 2.6 \\
    \quad $+f(0)$ & +8.5 & +7.0 & +9.3 & +8.1 & +8.2 & 0.4 \\
    \midrule
    15.3$\times$ no mix & +1.2 & -1.6 & +3.0 & +2.1 & +1.2 & 6.6 \\
    \quad $+f(0)$ & +7.3 & +5.8 & +8.8 & +7.2 & +7.3 & 0.5 \\
    15.3$\times$ +mix & +5.4 & +4.7 & +6.9 & +3.9 & +5.2 & 2.5 \\
    \quad $+f(0)$ & +7.3 & +6.1 & +8.5 & +7.2 & +7.3 & 0.4 \\
    \midrule
    M2L +mix & -9.9 & -12.3 & -9.9 & -12.7 & -11.2 & 8.4 \\
\bottomrule
\end{tabular}

  \caption{Oracle stem subtraction on MUSDB18 (SI-SDR, dB; 5{,}880 units,
  49 tracks). M2L denotes Music2Latent. ``$+f(0)$'' applies recentering to the same GAN checkpoint.
  ``gap'' is the shortfall from each model's reconstruction SI-SDR for the
  true residual (the linearity tax; lower is better).}
  \label{tab:subtraction}
\end{table}

\tabref{tab:subtraction} shows that raw subtraction favours mixing
supervision: $\Ldec$ improves SI-SDR from
$3.4$ to $5.8$\,dB at $7.66\times$, and the combined recipe improves it
from $1.2$ to $5.2$\,dB at $15.3\times$. Recentering improves all five GAN
configurations by $2.0$--$6.0$\,dB and brings their aggregate scores within
about $0.5$\,dB of their own reconstruction references, although individual
tracks have larger gaps. After correction, the two treatment--control
differences are $-0.03$\,dB $[-0.15,0.10]$ and $0.01$\,dB
$[-0.19,0.21]$, respectively.

On the independent MoisesDB corpus, raw advantages of $2.3$--$4.5$\,dB
shrink by $87$--$95$\% after recentering. The absolute between-model gap
shrinks on $229$--$236$ of the 240 recordings. A smaller advantage remains:
$0.29$\,dB for $\Ldec$, $0.12$\,dB for $\Ldec{+}$disc at $7.66\times$,
and $0.42$\,dB for the combined recipe at $15.3\times$, each with an
interval excluding zero. The held-out evaluation therefore supports a large
origin effect while retaining a small training-associated difference.
The registered activity filter keeps units whose removed stem is at least
$-50$\,dBFS and within $30$\,dB of the mixture energy. It retains
$24{,}648$ of $28{,}800$ units; vocals have the lowest retention ($70.4\%$).
This filter tests activity of the removed stem, so it does not guarantee
that the remaining mixture used as the SI-SDR reference is nonsilent.

On MUSDB18, the unchanged mixture is an essential baseline: it scores $17.1$\,dB
against the residual, above every model's reconstruction reference, because
the removed stem often contributes little energy. High residual SI-SDR alone
therefore does not establish successful removal. Relative to a second
baseline, decoding the mixture and stem separately before waveform
subtraction, recentered latent subtraction gains $0.7$--$0.9$\,dB; raw
latent subtraction performs worse.

\subsection{Compression and architecture}\label{sec:scope}
At $15.3\times$ compression, the combined recipe improves original-test
mixing SI-SDR from $6.8$ to $8.6$\,dB and MixRate from $1.19$ to $1.05$.
Reconstruction SI-SDR is similar ($10.0$ versus $9.9$\,dB), while FAD
increases from $0.052$ to $0.061$. This comparison concerns two separately
initialized runs of the full recipes. Mixing gains occur in all three test
domains at both compression rates ($1.2$--$3.2$\,dB). At $7.66\times$, the
gain diagnostic also shows less level error: decoded mixtures are $1.6$\,dB too quiet
without mixing supervision and $0.9$\,dB too quiet with it.

In Music2Latent, fine-tuning with mixed-latent supervision under the native
consistency loss improves ground-truth mixing SI-SDR from $-9.5$ to
$-7.8$\,dB relative to a continued-training control. Continued training
alone changes the mixing score by less than $0.2$\,dB. Absolute waveform
agreement remains poor in this experiment. Because the systems differ in
architecture and training, this result does not isolate the reason for that
performance or establish a limitation of consistency autoencoders generally.

\section{Conclusion}\label{sec:conclusion}
Mixing supervision improves waveform mixing in the evaluated GAN
autoencoders, including on independent audio. Decoder agreement can improve
as waveform fidelity decreases, while recentering removes most of the raw
subtraction advantage without changing reconstruction or convex mixing.
Ground-truth mixing fidelity, decoder agreement, and origin-corrected
subtraction should therefore be reported separately.

The comparisons do not isolate the contributions of latent interpolation,
mixture targets, and the extra waveform loss. Each configuration has one
training run, and bootstrap intervals cover evaluation variation for fixed
checkpoints. Fine-tunes share initial weights; from-scratch runs do not.
Checkpoint selection used the original test split. Evaluation is limited to
mono audio, oracle subtraction, and objective metrics without a listening study.

\clearpage
\bibliographystyle{IEEEtran}
\bibliography{references}

\end{document}
